%% 5-17 近くでは多項式で近似できる：テイラー展開（第11回：ページ制限なし。なぜその数になるかを書き足した）
\newcommand\bkvqmidashi[1]{\par\Needspace{6\baselineskip}\noindent\textbf{#1}\par}
\bkchapter{5}{17}{近くでは多項式で近似できる：\\テイラー展開}{\bklevel{4}}{近くでは多項式で近似できる：テイラー展開}
\bkmichishirube{5-8,5-14,5-15}{6-6,6-10,9-1}

\begin{goal}
\begin{itemize}[label={\color{bkTheme}●}, leftmargin=1.6\zw, itemsep=0pt, topsep=1pt]
\item \bklevel{3}\enspace 1次近似（接線）に2次の項をたすと、近似がよくなることを説明できる。
\item \bklevel{4}\enspace $e^x$、$\sin x$、$\cos x$ のテイラー展開を作り、近似値を計算できる。
\item \bklevel{4}\enspace $\sin\theta\approx\theta$ を使って、振り子の周期を求められる。
\end{itemize}
\end{goal}

\begin{hatena}
\noindent\begin{minipage}[c]{0.6\linewidth}
電卓で $\sin0.5$ を計算すると、0.4794255… が一瞬で出てきます。

でも、電卓やコンピュータの中の回路がもともとできるのは、たし算・ひき算・かけ算・わり算のような計算と、そのくり返しだけです。円もグラフも書けない計算機は、どうやって $\sin$ の値を出しているのでしょう。
\end{minipage}\hfill
\begin{minipage}[c]{0.36\linewidth}
\centering
\begin{tikzpicture}[x=1mm, y=1mm]
  \draw[bkInk!70, line width=1pt, rounded corners=3pt, fill=bkPaper] (0,0) rectangle (30,44);
  \draw[bkInk!50, line width=0.6pt, fill=white] (3,33) rectangle (27,41);
  \node[font=\ttfamily\small, anchor=east] at (26,37) {0.4794255};
  \foreach \i in {0,1,2,3} \foreach \j in {0,1,2} {\draw[bkInk!40, fill=bkTheme!10, rounded corners=1pt] (3+\i*6.5,3+\j*7) rectangle ++(5,5);}
  \foreach \i in {2,3} {\draw[bkInk!40, fill=bkTheme!10, rounded corners=1pt] (3+\i*6.5,24) rectangle ++(5,5);}
  \draw[bkInk!40, fill=bkTheme!10, rounded corners=1pt] (3,24) rectangle (14.5,29);
  \node[font=\sffamily\footnotesize] at (8.75,26.5) {sin};
  \node[font=\sffamily\bfseries\small] at (38,22) {？};
\end{tikzpicture}
\end{minipage}
\end{hatena}

答えのかぎは、「なめらかな関数は、ある点の近くでは、$x$ と $x^2$ と $x^3$ と……の、たし算の式（多項式）で表せる」という考えです。たし算とかけ算の式なら、計算機でも計算できます。

\clearpage
%% ---------------- やってみよう・思い出そう ----------------
\begin{yattemiyou}[表計算で、$e^{0.5}$ に近づけよう]
$e^{0.5}$ の正しい値は 1.6487212…。次の式に $x=0.5$ を入れて、項を1つずつたしていきます。$3!$ は $1\times2\times3=6$、$4!$ は $1\times2\times3\times4=24$、$5!$ は $1\times2\times3\times4\times5=120$ のことです（くわしくは③）。
\[ 1+x+\frac{x^2}{2}+\frac{x^3}{3!}+\frac{x^4}{4!}+\frac{x^5}{5!}+\cdots \]
\begin{bkcenter}
\bksheet{3}{28}{%
1 & たした項 & 項の値 & ここまでの合計 \\
2 & 1 & 1 & 1 \\
3 & $x$ & 0.5 & 1.5 \\
4 & $\frac{x^2}{2}$ & 0.125 & 1.625 \\
5 & $\frac{x^3}{6}$ & 0.020833 & 1.645833 \\
6 & $\frac{x^4}{24}$ & 0.002604 & 1.648438 \\
7 & $\frac{x^5}{120}$ & 0.000260 & 1.648698 \\}
\bktabcaption{項をたすたびに、正しい値 1.648721… に近づく。}
\end{bkcenter}
\end{yattemiyou}

\noindent たす項は、どんどん小さくなり、合計は正しい値にどんどん近づきます。かけ算とたし算だけで、$e^{0.5}$ が計算できました。この式はどこから来たのでしょう。なぜ2、6、24、120でわるのでしょう。

\begin{omoidasou}
\bkrecall{5-4}{$(x^n)'=nx^{n-1}$。定数の微分は0。定数倍は前に出せる。}
\bkrecall{5-7}{2回微分 $f''(x)$ は、グラフの曲がり方を表す。$f''>0$ なら下に凸（上に曲がる）。}
\bkrecall{5-8}{1次近似：$x$ が $a$ に近いとき $f(x)\approx f(a)+f'(a)(x-a)$（接線で近似）}
\bkrecall{5-14}{$(e^x)'=e^x$、$e^0=1$。$(x^r)'=rx^{r-1}$（$r$ が分数でもよい）。$(\ln x)'=\dfrac{1}{x}$}
\bkrecall{5-15}{$\sin$ と $\cos$ は、微分するたびに $\sin\to\cos\to-\sin\to-\cos\to\sin$ と回る。}
\end{omoidasou}

\clearpage
%% ---------------- 数学の道具 ----------------
\bksection{doguu}{数学の道具}

\bksubsection{1}{接線で近似する（1次近似）}

\label{bk:5-17:ichiji}5-8では、グラフを接線で近似しました。$x=0$ の近くでは、接線は「点 $(0,\ f(0))$ を通り、傾き $f'(0)$ の直線」なので、
\[ f(x)\approx f(0)+f'(0)\,x \]
$e^x$ では、$f(0)=e^0=1$。$(e^x)'=e^x$（5-14）なので $f'(0)=e^0=1$。だから $e^x\approx1+x$ です。どのくらい合うか、表で見てみましょう。
\begin{bkcenter}
\begin{bktable}{colspec={Q[l,m,wd=40mm] *{3}{Q[c,m,wd=20mm]}}, rowsep=1.5pt}
$x$ & 0.1 & 0.5 & 1 \\
正しい値 $e^x$ & 1.10517 & 1.64872 & 2.71828 \\
1次近似 $1+x$ & 1.1 & 1.5 & 2 \\
ずれ & 0.005 & 0.149 & 0.718 \\
\end{bktable}
\end{bkcenter}
$x=0.1$ ではよく合いますが、$x$ が0からはなれるほど、ずれが大きくなります。接線はまっすぐなのに、$e^x$ のグラフは上に曲がっている（下に凸、5-7）ので、だんだんはなれていってしまうのです。

\bksubsection{2}{曲がり方も合わせる（2次近似）}

\label{bk:5-17:niji}\bkidx{にじきんじ}{2次近似}そこで、まっすぐな直線のかわりに、曲がった放物線を使います。放物線 $p(x)=a_0+a_1x+a_2x^2$ には、自由に決められる数が3つ（$a_0$、$a_1$、$a_2$）あるので、$x=0$ で次の3つを $f(x)$ に合わせられます。
\begin{bkcenter}
\begin{bktable}{colspec={Q[l,m,wd=56mm] Q[l,m,wd=56mm]}, rowsep=1.5pt}
合わせるもの & 表すもの \\
値 $p(0)=f(0)$ & 同じ点を通る \\
1回微分 $p'(0)=f'(0)$ & 傾きが同じ（接線が同じ） \\
2回微分 $p''(0)=f''(0)$ & 曲がり方が同じ \\
\end{bktable}
\end{bkcenter}
$p(x)$ を微分していきます（5-4）。
\[ p(x)=a_0+a_1x+a_2x^2,\qquad p'(x)=a_1+2a_2x,\qquad p''(x)=2a_2 \]
それぞれに $x=0$ を入れると、$x$ のつく項は0になって、
\[ p(0)=a_0,\qquad p'(0)=a_1,\qquad p''(0)=2a_2 \]
これを $f(0)$、$f'(0)$、$f''(0)$ と等しくすると、$a_0=f(0)$、$a_1=f'(0)$、そして $2a_2=f''(0)$ から $a_2=\dfrac{f''(0)}{2}$。
\[ f(x)\approx f(0)+f'(0)\,x+\frac{f''(0)}{2}\,x^2 \]
$x^2$ の項で2でわるのは、$x^2$ を2回微分すると2が出てくるので、それを打ち消すためです。

$e^x$ は何回微分しても $e^x$ なので $f''(0)=e^0=1$。だから $e^x\approx1+x+\dfrac{x^2}{2}$ です。
\begin{bkcenter}
\begin{bktable}{colspec={Q[l,m,wd=40mm] *{3}{Q[c,m,wd=20mm]}}, rowsep=1.5pt}
$x$ & 0.1 & 0.5 & 1 \\
正しい値 $e^x$ & 1.10517 & 1.64872 & 2.71828 \\
2次近似 $1+x+\frac{x^2}{2}$ & 1.105 & 1.625 & 2.5 \\
ずれ & 0.0002 & 0.024 & 0.218 \\
\end{bktable}
\end{bkcenter}
1次近似より、ずれがずっと小さくなりました。たとえば $x=0.5$ では、ずれが0.149から0.024にへりました。

\begin{bkfigure}
\begin{tikzpicture}
\begin{axis}[bkplot, scale only axis, width=84mm, height=40mm, xmin=-2, xmax=2, ymin=-0.6, ymax=6.2, xtick={-2,-1,1,2}, ytick={2,4,6}, ylabel={$y$},
    axis x line=bottom, x axis line style={draw=none}, xticklabel style={yshift=-0.5mm},
    yticklabel style={fill=white, inner sep=1pt, outer sep=2.3pt}]
  \draw[bkInk, line width=0.8pt, -{Stealth[length=5pt]}] (axis cs:-2,0) -- (axis cs:2.12,0) node[anchor=west, font=\sffamily\small, text=bkInk] {$x$};
  \addplot[bkcurve, domain=-2:1.82, samples=80, line width=2.2pt] {exp(x)};
  \addplot[bkthin, domain=-1.6:2] {1+x};
  \addplot[bkcurve2, domain=-2:2, samples=60] {1+x+x^2/2};
  \addplot[bkcurve3, domain=-2:1.95, samples=60] {1+x+x^2/2+x^3/6};
  \node[font=\sffamily\footnotesize, text=bkBlue, anchor=east] at (axis cs:1.6,6.35) {$e^x$};
  \node[font=\sffamily\footnotesize, text=bkGreen!70!black, anchor=west, fill=white, inner sep=1pt] at (axis cs:2.05,6.1) {$1+x+\frac{x^2}{2}+\frac{x^3}{6}$};
  \node[font=\sffamily\footnotesize, text=bkOrange!80!black, anchor=west, fill=white, inner sep=1pt] at (axis cs:2.05,5) {$1+x+\frac{x^2}{2}$};
  \node[font=\sffamily\footnotesize, text=bkGray, anchor=west, fill=white, inner sep=1pt] at (axis cs:2.05,3) {$1+x$};
\end{axis}
\end{tikzpicture}
\bkfigcaption{$e^x$（太い青）と、その近似。項をふやすほど、$x=0$ から遠くまでよく合う。}
\end{bkfigure}

\bksubsection{3}{何回でも合わせる：階乗とテイラー展開}

\label{bk:5-17:kaijo}同じように $x^3$、$x^4$、… の項をたして、3回微分、4回微分、… まで合わせていきます。まず、$x^3$ を3回微分してみましょう。
\[ x^3\ \xrightarrow{\text{微分}}\ 3x^2\ \xrightarrow{\text{微分}}\ 3\times2x\ \xrightarrow{\text{微分}}\ 3\times2\times1=6 \]
$x^4$ を4回微分すると、$4\times3\times2\times1=24$。指数が前に出て、1ずつへりながら、かけ算がつながっていきます。一般に、$x^n$ を $n$ 回微分すると、1から $n$ までの整数をすべてかけた数になります。これを $n!$ と書き、「$n$ の\textbf{階乗}」と読みます。$0!=1$ と決めておきます。\bkidx{かいじょう}{階乗}
\begin{bkcenter}
\begin{tblr}{colspec={Q[c,m]*{7}{Q[c,m]}}, row{1}={bg=bkTheme!12, font=\sffamily\small}, hline{1,2,Z}={0.6pt,bkGray}, colsep=7pt}
$n$ & 0 & 1 & 2 & 3 & 4 & 5 & 6 \\
$n!$ & 1 & 1 & 2 & 6 & 24 & 120 & 720 \\
\end{tblr}
\end{bkcenter}
表の右へ1つ進むと、$n$ をかけることになります。たとえば $5!=4!\times5$ なので、$24\times5=120$ です。

\bkvqmidashi{3次の多項式で確かめる}
$p(x)=a_0+a_1x+a_2x^2+a_3x^3$ を3回微分し、$x=0$ を入れます。
\begin{bkcenter}
\begin{bktable}{colspec={Q[l,m,wd=12mm] Q[l,m,wd=62mm] Q[c,m,wd=26mm]}, rowsep=1.5pt}
 & 式 & $x=0$ を入れると \\
$p$ & $a_0+a_1x+a_2x^2+a_3x^3$ & $a_0$ \\
$p'$ & $a_1+2a_2x+3a_3x^2$ & $a_1$ \\
$p''$ & $2a_2+3\times2\,a_3x$ & $2a_2$ \\
$p'''$ & $3\times2\times1\,a_3$ & $6a_3=3!\,a_3$ \\
\end{bktable}
\end{bkcenter}
3回微分すると、$x^2$ までの項は消えて（定数の微分は0）、$x^3$ の項だけが $3!\,a_3$ として残ります。だから3回微分を $f'''(0)$ に合わせるには $a_3=\dfrac{f'''(0)}{3!}$ にします。もっと次数の高い項があっても、$x=0$ を入れたとき $x$ が残って0になるので、じゃまをしません。

同じように、$a_nx^n$ の項を $n$ 回微分すると $a_n\times n!$ なので、$n$ 回微分を合わせるには $a_n=\dfrac{f^{(n)}(0)}{n!}$ にします（$f^{(n)}$ は $f$ を $n$ 回微分したもの）。
\[ f(x)=f(0)+f'(0)\,x+\frac{f''(0)}{2!}\,x^2+\frac{f'''(0)}{3!}\,x^3+\cdots+\frac{f^{(n)}(0)}{n!}\,x^n+\cdots \]
これを $f(x)$ の\textbf{テイラー展開}（$x=0$ のまわりの展開）といいます。項をふやすほど、$x=0$ から遠くまで、よく合うようになります。\bkidx{ていらーてんかい}{テイラー展開}

\bksubsection{4}{$e^x$ のテイラー展開}

\label{bk:5-17:e}$e^x$ は何回微分しても $e^x$ で、$x=0$ ではいつも $e^0=1$。だから、③の式の $f(0)$、$f'(0)$、$f''(0)$、… が全部1になって、
\[ e^x=1+x+\frac{x^2}{2!}+\frac{x^3}{3!}+\frac{x^4}{4!}+\cdots \]
やってみようの式は、これでした。$x=1$ を入れると、$e^1=e$ が計算できます。
\begin{bkcenter}
\begin{bktable}{colspec={Q[l,m,wd=21mm] *{8}{Q[c,m,wd=12mm]}}, rowsep=1.5pt, colsep=1pt, cells={font=\footnotesize}}
たす項 & 1 & 1 & $\frac{1}{2}$ & $\frac{1}{6}$ & $\frac{1}{24}$ & $\frac{1}{120}$ & $\frac{1}{720}$ & $\frac{1}{5040}$ \\
ここまでの合計 & 1 & 2 & 2.5 & 2.6667 & 2.7083 & 2.7167 & 2.71806 & 2.71825 \\
\end{bktable}
\bktabcaption{$e=2.718281\cdots$ に近づいていく。}
\end{bkcenter}
5-14の $\Bigl(1+\dfrac{1}{n}\Bigr)^n$ では、$n=10000$ にしても2.71814…で、小数3けたまでしか合いませんでした。こちらは、わずか8項のたし算で、小数4けたまで合います。

\bksubsection{5}{$\sin x$ と $\cos x$ のテイラー展開}

\label{bk:5-17:sin}$\sin x$ を微分していくと $\cos x$、$-\sin x$、$-\cos x$、$\sin x$、… と回ります（5-15）。$\sin0=0$、$\cos0=1$ を使って、$x=0$ での値を表にします。
\begin{bkcenter}
\begin{tblr}{colspec={Q[l,m]*{6}{Q[c,m]}}, row{1}={bg=bkTheme!12, font=\sffamily\small}, hline{1,2,Z}={0.6pt,bkGray}, colsep=7pt}
 & $f$ & $f'$ & $f''$ & $f'''$ & $f^{(4)}$ & $f^{(5)}$ \\
$\sin x$ & 0 & 1 & 0 & $-1$ & 0 & 1 \\
$\cos x$ & 1 & 0 & $-1$ & 0 & 1 & 0 \\
\end{tblr}
\end{bkcenter}
たとえば $\sin x$ の3回微分は $-\cos x$ なので、$x=0$ で $-\cos0=-1$ です。③の式に入れると、0の所の項は消えるので、
\[ \sin x=x-\frac{x^3}{3!}+\frac{x^5}{5!}-\frac{x^7}{7!}+\cdots,\qquad \cos x=1-\frac{x^2}{2!}+\frac{x^4}{4!}-\frac{x^6}{6!}+\cdots \]
$\sin x$ は奇数の次数だけ、$\cos x$ は偶数の次数だけで、$+$ と $-$ が交代します。

\bkvqmidashi{確かめ}
$\sin x$ の展開を、項ごとに微分してみます。たとえば $\dfrac{x^3}{3!}$ を微分すると、
\[ \frac{3x^2}{3!}=\frac{3x^2}{1\times2\times3}=\frac{x^2}{1\times2}=\frac{x^2}{2!} \]
のように、階乗が1つへります。
\[ \Bigl(x-\frac{x^3}{3!}+\frac{x^5}{5!}-\cdots\Bigr)'=1-\frac{3x^2}{3!}+\frac{5x^4}{5!}-\cdots=1-\frac{x^2}{2!}+\frac{x^4}{4!}-\cdots \]
ちゃんと $\cos x$ の展開になります（$(\sin x)'=\cos x$ と合う）。\bkidx{じすう}{次数}

\bkvqmidashi{\bkhatenaref の答え}
$\sin x$ の展開のはじめの3項に $x=0.5$ を入れます。$0.5^3=0.125$、$0.5^5=0.03125$ です。
\begin{align*}
\sin0.5&\approx0.5-\frac{0.5^3}{3!}+\frac{0.5^5}{5!}=0.5-\frac{0.125}{6}+\frac{0.03125}{120}\\
&=0.5-0.0208333+0.0002604=0.4794271
\end{align*}
3項で、電卓の0.4794255と小数5けたまで合います。たし算・ひき算・かけ算・わり算だけで、$\sin$ の値が出せました。コンピュータの中の $\sin$ や $\exp$ も、このような多項式の考えで計算しています（実際には、計算を速く正確にするための工夫をさらに加えています。電卓には、たし算と「けたをずらす」計算だけで $\sin$ を求める別の方法を使うものもあります）。

\begin{bkfigure}
\begin{tikzpicture}
\begin{axis}[bkplot, scale only axis, width=84mm, height=38mm, xmin=-3.6, xmax=3.6, ymin=-2, ymax=2, xtick={-3.1416,-1.5708,1.5708,3.1416}, xticklabels={$-\pi$,$-\dfrac{\pi}{2}$,$\dfrac{\pi}{2}$,$\pi$}, ytick={-1,1},
    axis x line=bottom, x axis line style={draw=none}, xticklabel style={anchor=base, yshift=-5.5mm},
    yticklabel style={fill=white, inner sep=1pt, outer sep=2.3pt}]
  \draw[bkInk, line width=0.8pt, -{Stealth[length=5pt]}] (axis cs:-3.6,0) -- (axis cs:3.75,0) node[anchor=west, font=\sffamily\small, text=bkInk] {$x$};
  \addplot[bkcurve, domain=-3.6:3.6, samples=120, line width=2.2pt] {sin(deg(x))};
  \addplot[bkthin, domain=-2:2] {x};
  \addplot[bkcurve2, domain=-3.0:3.0, samples=80] {x-x^3/6};
  \addplot[bkcurve3, domain=-3.6:3.6, samples=80] {x-x^3/6+x^5/120};
  \node[font=\sffamily\footnotesize, text=bkBlue, anchor=east] at (axis cs:-3.65,0.44) {$\sin x$};
  \node[font=\sffamily\footnotesize, text=bkGray, anchor=west, fill=white, inner sep=1pt] at (axis cs:2.05,2) {$x$};
  \node[font=\sffamily\footnotesize, text=bkOrange!80!black, anchor=west, fill=white, inner sep=1pt] at (axis cs:3.07,-1.5) {$x-\frac{x^3}{6}$};
  \node[font=\sffamily\footnotesize, text=bkGreen!70!black, anchor=west, fill=white, inner sep=1pt] at (axis cs:3.67,0.86) {$x-\frac{x^3}{6}+\frac{x^5}{120}$};
\end{axis}
\end{tikzpicture}
\bkfigcaption{$\sin x$（太い青）と、そのテイラー展開を1項、2項、3項で止めたもの。3項なら、$-\frac{\pi}{2}$ から $\frac{\pi}{2}$ くらいまで、ほとんど重なる。}
\end{bkfigure}

\bksubsection{6}{小さい $x$ では、はじめの項だけでよい}

\label{bk:5-17:chiisai}$x$ が小さいと、$x^2$、$x^3$ はもっと小さくなります。たとえば $x=0.1$ なら $x^2=0.01$、$x^3=0.001$。そこで、はじめの1つか2つの項だけを使う近似がよく使われます。
\begin{bkcenter}
\begin{tblr}{colspec={Q[l,m]Q[l,m]Q[l,m]}, row{1}={bg=bkTheme!12, font=\sffamily\small}, hline{1,2,Z}={0.6pt,bkGray}, hline{3,4,5}={0.3pt,bkGray}, colsep=5.5pt}
近似の式 & 次の項（ずれの目安） & 使われる所 \\
$\sin\theta\approx\theta$ & $-\dfrac{\theta^3}{6}$ & 振り子、光の屈折、小さな傾き \\
$\cos\theta\approx1-\dfrac{\theta^2}{2}$ & $+\dfrac{\theta^4}{24}$ & 振り子の高さ、アンテナの向き \\
$e^x\approx1+x$ & $+\dfrac{x^2}{2}$ & 小さな利率、温度による変化 \\
$(1+x)^r\approx1+rx$ & $+\dfrac{r(r-1)}{2}x^2$ & 熱による伸び、小さなずれの計算 \\
\end{tblr}
\end{bkcenter}

\bkvqmidashi{$(1+x)^r$ の近似の作り方}
$f(x)=(1+x)^r$ を微分します。内側 $1+x$ の微分は1なので、合成関数の微分（5-6）と $(x^r)'=rx^{r-1}$（5-14）から、
\[ f'(x)=r(1+x)^{r-1},\qquad f''(x)=r(r-1)(1+x)^{r-2} \]
$x=0$ を入れると、$1+0=1$ は何乗しても1なので、$f(0)=1$、$f'(0)=r$、$f''(0)=r(r-1)$。②の式に入れると、
\[ (1+x)^r\approx1+rx+\frac{r(r-1)}{2}x^2 \]
はじめの2項が $1+rx$、3つめの項が「次の項」です。

\bkvqmidashi{$\sin\theta\approx\theta$ はどのくらい合う？}
表で確かめましょう（$\theta$ はラジアン。度に $\frac{\pi}{180}$ をかけて求める）。
\begin{bkcenter}
\begin{tblr}{colspec={Q[c,m]*{5}{Q[c,m]}}, row{1}={bg=bkTheme!12, font=\sffamily\small}, hline{1,2,Z}={0.6pt,bkGray}, colsep=7pt}
角（度） & 5° & 10° & 20° & 30° & 45° \\
$\theta$（rad） & 0.0873 & 0.1745 & 0.3491 & 0.5236 & 0.7854 \\
$\sin\theta$ & 0.0872 & 0.1736 & 0.3420 & 0.5 & 0.7071 \\
ずれ（\%） & 0.13 & 0.51 & 2.1 & 4.7 & 11 \\
\end{tblr}
\end{bkcenter}
ずれの割合は、次の項の大きさ $\dfrac{\theta^3}{6}$ を、近似の値 $\theta$ でわった $\dfrac{\theta^3}{6}\div\theta=\dfrac{\theta^2}{6}$ くらいです。10°（0.1745 rad）なら $\dfrac{0.1745^2}{6}=0.0051$ で0.5\,\%、20°で2\,\%。表と合います。「どこまで近似してよいか」は、次の項の大きさで見積もります。

\bksubsection{7}{$x=a$ のまわりの展開}

\label{bk:5-17:a}$\sqrt{4.1}$ のように、0から遠い所の値を知りたいときは、近くの計算しやすい点 $a$（ここでは4）のまわりで展開します。$a$ からのずれ $x-a$ を新しい小さな数と見て、$x-a$ の多項式で近似するのです。考え方は③と同じで、$x=a$ で値・傾き・曲がり方…を合わせると、
\[ f(x)=f(a)+f'(a)(x-a)+\frac{f''(a)}{2!}(x-a)^2+\frac{f'''(a)}{3!}(x-a)^3+\cdots \]
$a=0$ にすると、③の式にもどります。

\bkvqmidashi{例：$\sqrt{4.1}$ を求める}
$f(x)=\sqrt{x}=x^{\frac{1}{2}}$、$a=4$ とします（$\sqrt4=2$ とすぐ分かるから）。$x-a=4.1-4=0.1$ です。
\begin{align*}
f'(x)&=\frac{1}{2}x^{-\frac{1}{2}}=\frac{1}{2\sqrt{x}}, & f'(4)&=\frac{1}{2\times2}=\frac{1}{4}\\
f''(x)&=\frac{1}{2}\times\Bigl(-\frac{1}{2}\Bigr)x^{-\frac{3}{2}}=-\frac{1}{4\,x\sqrt{x}}, & f''(4)&=-\frac{1}{4\times4\times2}=-\frac{1}{32}
\end{align*}
\[ \sqrt{4.1}\approx2+\frac{1}{4}\times0.1-\frac{1}{32}\times\frac{0.1^2}{2}=2+0.025-0.00015625=2.02484 \]
電卓の値は2.0248457。2次までで、小数5けた合いました。

\Needspace{12\baselineskip}
\bksubsection{8}{振り子は、小さくゆれれば単振動}

\label{bk:5-17:furiko}長さ $L$ の糸につるした、質量 $m$ のおもりの振り子です。

\noindent\begin{minipage}[c]{0.4\linewidth}
\centering
\begin{tikzpicture}[x=1mm, y=1mm]
  \fill[bkGray!40] (-14,0) rectangle (14,2);
  \draw[bkGray, dashed, line width=0.5pt] (0,0) -- (0,-36.5);
  \def\th{25}
  \draw[bkInk, line width=0.9pt] (0,0) -- (-90+\th:36);
  \draw[bkInk, line width=0.5pt] (0,-9) arc (-90:-90+\th:9);
  \node[font=\sffamily\scriptsize] at (-90+\th/2:12) {$\theta$};
  \node[font=\sffamily\scriptsize, anchor=west, xshift=1mm] at (-90+\th:18) {$L$};
  \draw[bkBlue!50, dashed, line width=0.6pt] (-90-\th:36) arc (-90-\th:-90+\th:36);
  \fill[bkOrange] (-90+\th:36) circle (2.6);
  \draw[bk arrow, bkInk, line width=1pt] (-90+\th:36) -- ++(0,-14) node[anchor=north, font=\sffamily\scriptsize] {$mg$};
  \draw[bk arrow, bkRed, line width=1.2pt] (-90+\th:36) -- ++(180+\th:8);
  \node[anchor=north east, font=\sffamily\scriptsize, text=bkRed, inner sep=0pt] at (9,-38.2) {$mg\sin\theta$};
\end{tikzpicture}
\end{minipage}\hfill
\begin{minipage}[c]{0.56\linewidth}
おもりにはたらく重力は $mg$（質量×重力加速度。$g=\qty{9.8}{m/s^2}$）で、真下を向いています。この力を、糸の向きと、円の弧にそった向きの2つに分けます（4-6の「斜面の重力を2つに分ける」と同じ）。

糸と真下の向きのあいだの角が $\theta$ なので、弧にそった向きの成分は $mg\sin\theta$（赤い矢印）です。糸の向きの成分は、糸が引く力といっしょに円の中心の向きにはたらくだけで、おもりを弧にそって進めたりもどしたりはしません。弧にそってふれを引きもどすのは、この $mg\sin\theta$ だけです。
\end{minipage}\par\vspace{3pt}
弧にそってはかった位置は $s=L\theta$（弧の長さ＝半径×角、4-2）です。$L$ は定数なので、2回微分すると、弧にそった加速度は $\ddot{s}=L\ddot{\theta}$。力＝質量×加速度（5-5）から、
\[ mL\ddot{\theta}=-mg\sin\theta \]
（$-$ は、ふれと反対向きに引きもどすから。）両辺を $mL$ でわると、
\[ \ddot{\theta}=-\frac{g}{L}\sin\theta \]
質量 $m$ は両辺から消えました。このままでは、5-15の単振動の形 $\ddot{x}=-\omega^2x$（右が $x$ そのものの何倍か）ではありません。でも、ふれが小さければ $\sin\theta\approx\theta$（⑥）なので、
\[ \ddot{\theta}\approx-\frac{g}{L}\,\theta \]
これは $\omega^2=\dfrac{g}{L}$、つまり $\omega=\sqrt{\dfrac{g}{L}}$ の単振動です。周期は5-15の $T=\dfrac{2\pi}{\omega}$ から、
\[ T=\frac{2\pi}{\omega}=2\pi\div\sqrt{\frac{g}{L}}=2\pi\sqrt{\frac{L}{g}} \]
（$\sqrt{\frac{g}{L}}$ でわることは、$\sqrt{\frac{L}{g}}$ をかけることと同じ。）この式には、ふれはば $\theta$ も質量 $m$ も入っていません。\textbf{小さいふれの周期は、ふれはばにもおもりの質量にもよらず、糸の長さだけで決まります}。2-8で「振り子の周期は糸の長さの $\sqrt{\phantom{x}}$ に比例する」と表で調べたことの理由が、ここで分かりました。たとえば $L=\qty{1}{m}$ なら $T=2\times3.14\times\sqrt{\dfrac{1}{9.8}}\approx\qty{2.0}{s}$ です。\bkidx{ふりこのしゅうき}{振り子の周期}

ふれが大きくなると $\sin\theta\approx\theta$ がずれてきて、周期は少し長くなります。
\begin{bkfigure}
\begin{tikzpicture}
\begin{axis}[bkplot, width=110mm, height=34mm, xmin=0, xmax=95, ymin=0.98, ymax=1.22, xtick={0,10,20,30,45,60,90}, xticklabels={0°,10°,20°,30°,45°,60°,90°}, ytick={1,1.1,1.2}, yticklabels={1倍,1.1倍,1.2倍}, xlabel={ふれはば}]
  \addplot[bkcurve, mark=*, mark size=1.4pt] coordinates {(0,1) (10,1.0019) (20,1.0077) (30,1.0174) (40,1.0313) (50,1.0498) (60,1.0732) (70,1.1021) (80,1.1375) (90,1.1803)};
  % ラベルは、横の格子の線（1.1倍・1.2倍）のあいだに置く
  \node[font=\sffamily\footnotesize, anchor=west, fill=white, inner sep=1pt] at (axis cs:3,1.15) {10°では0.2\,\%長いだけ};
  \node[font=\sffamily\footnotesize, anchor=east, fill=white, inner sep=1pt] at (axis cs:76,1.172) {90°では18\,\%長い};
\end{axis}
\end{tikzpicture}
\bkfigcaption{ふれはばと周期（小さいふれの周期を1倍とし、$\sin\theta$ のまま計算した値）。}
\end{bkfigure}

\Needspace{12\baselineskip}
\begin{motto}[5]
\textbf{オイラーの公式への入口}\par
$e^x$ の展開の $x$ に、「2乗すると $-1$ になる数」$j$（6-8）をかけた $j\theta$ を入れます。$j^2=-1$、$j^3=j^2\times j=-j$、$j^4=(j^2)^2=1$ を使って、$j$ のつかない項とつく項に分けると、
\[ e^{j\theta}=\Bigl(1-\frac{\theta^2}{2!}+\frac{\theta^4}{4!}-\cdots\Bigr)+j\Bigl(\theta-\frac{\theta^3}{3!}+\frac{\theta^5}{5!}-\cdots\Bigr)=\cos\theta+j\sin\theta \]
かっこの中は、⑤の $\cos\theta$ と $\sin\theta$ の展開そのものです。指数関数と三角関数が、1つの式でつながります。これが6-10のオイラーの公式で、交流回路の計算の土台になります。
\end{motto}

\begin{senmonsho}[大学の教科書ふうの書き方]
\begin{bkhonbun}
$f(x)$ の $x=a$ のまわりのテイラー展開は $\displaystyle f(x)=\sum_{n=0}^{\infty}\frac{f^{(n)}(a)}{n!}(x-a)^n$ である。特に $|\theta|\ll1$ のとき $\sin\theta=\theta+O(\theta^3)$ であり、振り子の運動方程式は線形化される。\bkidx{せんけいか}{線形化}
\end{bkhonbun}
\bkyaku{$\sum$ は5-10のたし算の記号。$\infty$ は「どこまでも」。$|\theta|\ll1$ は「$\theta$ が1よりずっと小さい」。$O(\theta^3)$（オーダー）は「$\theta^3$ くらいの小さな量」。「線形化」は、$\sin\theta$ を $\theta$ のような1次の式におきかえて、解きやすくすること。}
\end{senmonsho}

\bkkigocards{\bkkigocard{$n!$}{エヌのかいじょう}{1から $n$ までの整数をすべてかけた数。$0!=1$}{factorial}{5-17}}
{\bkkigocard{$f^{(n)}(x)$}{エフ・エヌかいびぶん}{$f(x)$ を $n$ 回微分したもの}{$n$th derivative}{5-17}}

\begin{taisetsu}
$\displaystyle f(x)=f(0)+f'(0)\,x+\frac{f''(0)}{2!}x^2+\frac{f'''(0)}{3!}x^3+\cdots$\\
$\displaystyle e^x=1+x+\frac{x^2}{2!}+\cdots$\quad $\displaystyle\sin x=x-\frac{x^3}{3!}+\cdots$\quad $\displaystyle\cos x=1-\frac{x^2}{2!}+\cdots$\\
小さい $\theta$ では $\sin\theta\approx\theta$。振り子の周期 $T=2\pi\sqrt{L/g}$
\end{taisetsu}

\begin{kotoba}
\small
\bkword{階乗}{1から $n$ までの整数をすべてかけた数 $n!$。$x^n$ を $n$ 回微分すると $n!$。}
\bkword{テイラー展開}{関数を、ある点の近くで多項式（たし算とかけ算の式）で表すこと。}
\bkword{線形化}{$\sin\theta\approx\theta$ のように、1次の式で近似して、式を解きやすくすること。}
\end{kotoba}

%% ---------------- 例題 ----------------
\bksection{reidai}{例題でたしかめよう}

\begin{reidai}[階乗の計算]
次の値を求めましょう。\quad（1）$5!$\quad（2）$\dfrac{6!}{4!}$\quad（3）$\dfrac{1}{4!}$ を小数で\par
\kangaekata $n!$ は1から $n$ までの整数をかけたもの（③）。わり算は、同じかけ算を約分する。\par
\shiki（1）$1\times2\times3\times4\times5=120$\par
\hspace*{2.6\zw}（2）$\dfrac{6!}{4!}=\dfrac{1\times2\times3\times4\times5\times6}{1\times2\times3\times4}=5\times6=30$\par
\hspace*{2.6\zw}（3）$4!=24$ なので、$\dfrac{1}{24}=0.041666\cdots$\par
\kotae （1）120\quad（2）30\quad（3）0.041666…
\end{reidai}

\begin{jibun}
次の値を求めましょう。\quad（1）$6!$\quad（2）$\dfrac{7!}{5!}$\quad（3）$\dfrac{1}{6!}$ を小数で
\bkjans{（1）\nobreak 720\quad（2）\nobreak $6\times7=42$\quad（3）\nobreak 0.0013888…}
\end{jibun}

\begin{reidai}[2次近似を作る]
次の関数の、$x=0$ のまわりの2次近似を求めましょう。\quad（1）$e^{2x}$\quad（2）$(1+x)^{10}$（これを使って $1.01^{10}$ も）\quad（3）$\sqrt{1+x}$\par
\kangaekata $x=0$ での $f$、$f'$、$f''$ の値を求め、②の式 $f(0)+f'(0)x+\frac{f''(0)}{2}x^2$ に入れる。\par
\shiki（1）$f'=2e^{2x}$、$f''=4e^{2x}$ より $f(0)=1$、$f'(0)=2$、$f''(0)=4$。$1+2x+\dfrac{4}{2}x^2$\par
\hspace*{2.6\zw}（2）$f'=10(1+x)^9$、$f''=90(1+x)^8$ より $f(0)=1$、$f'(0)=10$、$f''(0)=90$\par
\hspace*{4.6\zw}$1+10x+\dfrac{90}{2}x^2$。$x=0.01$ で $1+0.1+0.0045$（電卓1.10462）\par
\hspace*{2.6\zw}（3）$(1+x)^{\frac{1}{2}}$ と見る。$f'=\dfrac{1}{2}(1+x)^{-\frac{1}{2}}$、$f''=\dfrac{1}{2}\times\Bigl(-\dfrac{1}{2}\Bigr)(1+x)^{-\frac{3}{2}}$\par
\hspace*{4.6\zw}$f(0)=1$、$f'(0)=\dfrac{1}{2}$、$f''(0)=-\dfrac{1}{4}$。$1+\dfrac{1}{2}x+\dfrac{-\frac{1}{4}}{2}x^2$\par
\kotae （1）$1+2x+2x^2$\quad（2）$1+10x+45x^2$、$1.01^{10}\approx1.1045$\quad（3）$1+\dfrac{x}{2}-\dfrac{x^2}{8}$
\end{reidai}
（3）で $x=0.1$ とすると、$\sqrt{1.1}\approx1+0.05-0.00125=1.04875$。電卓の値1.04881とよく合います。

\begin{jibun}
$e^{-x}$ と $(1+x)^{\frac{1}{3}}$ の、$x=0$ のまわりの2次近似を求めましょう。
\bkjans{$e^{-x}\approx1-x+\dfrac{x^2}{2}$。$(1+x)^{\frac{1}{3}}$：$f'(0)=\frac{1}{3}$、$f''(0)=\frac{1}{3}\times\bigl(-\frac{2}{3}\bigr)=-\frac{2}{9}$ より \mbox{$1+\dfrac{x}{3}-\dfrac{x^2}{9}$}}
\end{jibun}

\begin{reidai}[$\sin$・$\cos$ の値を計算する]
$\sin0.2$ をテイラー展開の2項で、$\cos0.2$ を3項で計算しましょう。\par
\kangaekata ⑤の $\sin x$、$\cos x$ の展開の、はじめの項に $x=0.2$ を入れる。$0.2^2=0.04$、$0.2^3=0.008$、$0.2^4=0.0016$。\par
\shiki $\sin0.2\approx0.2-\dfrac{0.2^3}{3!}=0.2-\dfrac{0.008}{6}$\par
\hspace*{4.6\zw}$=0.2-0.001333=0.198667$\par
\hspace*{2.6\zw}$\cos0.2\approx1-\dfrac{0.2^2}{2!}+\dfrac{0.2^4}{4!}=1-\dfrac{0.04}{2}+\dfrac{0.0016}{24}$\par
\hspace*{4.6\zw}$=1-0.02+0.0000667=0.980067$\par
\kotae $\sin0.2\approx0.198667$（電卓0.198669）、$\cos0.2\approx0.980067$（電卓0.980067）
\end{reidai}

\begin{jibun}
$\sin0.3$ を2項で、$\cos0.3$ を3項（$1-\frac{x^2}{2!}+\frac{x^4}{4!}$）で計算しましょう。
\bkjans{$\sin0.3\approx0.3-0.0045=0.2955$（電卓0.29552）。$\cos0.3\approx1-0.045+0.0003375=0.955338$（電卓0.955336）}
\end{jibun}

\begin{reidai}[$e^x$ の値を計算する]
$e^{0.1}$ をテイラー展開の3項で、$e^{-0.1}$ を4項で計算しましょう。\par
\kangaekata ④の $e^x$ の展開に $x=0.1$、$x=-0.1$ を入れる。$x$ がマイナスのときは、奇数乗の項がマイナスになる（$(-0.1)^3=-0.001$）。\par
\shiki $e^{0.1}\approx1+0.1+\dfrac{0.1^2}{2}=1+0.1+0.005=1.105$\par
\hspace*{2.6\zw}$e^{-0.1}\approx1+(-0.1)+\dfrac{(-0.1)^2}{2}+\dfrac{(-0.1)^3}{6}$\par
\hspace*{4.6\zw}$=1-0.1+0.005-0.000167=0.904833$\par
\kotae $e^{0.1}\approx1.105$（電卓1.10517）、$e^{-0.1}\approx0.904833$（電卓0.904837）
\end{reidai}

\begin{jibun}
$e^{-0.2}$ を、テイラー展開の4項で計算しましょう。
\bkjans{$1-0.2+0.02-0.001333=0.818667$（電卓0.818731）}
\end{jibun}

\begin{reidai}[振り子の周期]
長さ\qty{1}{m}の振り子の、小さなゆれの周期を求めましょう（$g=\qty{9.8}{m/s^2}$、$\pi=3.14$）。\par
\kangaekata ⑧の周期の式 $T=2\pi\sqrt{\frac{L}{g}}$ に、$L=1$、$g=9.8$ を入れる。\par
\shiki $\dfrac{L}{g}=\dfrac{1}{9.8}=0.102$、$\sqrt{0.102}=0.319$\par
\hspace*{2.6\zw}$T=2\times3.14\times0.319=2.00$\par
\kotae 約2.0秒
\end{reidai}

\begin{jibun}
周期1秒の振り子の、糸の長さは何cmでしょう（$g=\qty{9.8}{m/s^2}$、$\pi^2=9.87$）。
\bkjans{$T=2\pi\sqrt{\dfrac{L}{g}}$ の両辺を2乗して $L=g\Bigl(\dfrac{T}{2\pi}\Bigr)^2=\dfrac{9.8\times1^2}{4\times9.87}\approx\qty{0.248}{m}$。約25\,cm}
\end{jibun}

\begin{reidai}[近似のずれを見積もる]
15°（0.2618 rad）での $\sin\theta\approx\theta$ のずれの割合を、次の項から見積もりましょう。\par
\kangaekata ずれの割合は、次の項 $\frac{\theta^3}{6}$ を $\theta$ でわった $\frac{\theta^2}{6}$（⑥）。\par
\shiki ずれの割合 $\approx\dfrac{\theta^2}{6}=\dfrac{0.2618^2}{6}=\dfrac{0.0685}{6}=0.0114$\par
\kotae 約1.1\,\%（電卓で計算した $\sin15^\circ=0.2588$ とくらべても1.2\,\%）
\end{reidai}

\begin{jibun}
$\cos\theta\approx1$ と近似したときのずれ（次の項 $\frac{\theta^2}{2}$）が0.01以下になるのは、$\theta$ が何radまでですか。
\bkjans{$\dfrac{\theta^2}{2}\leqq0.01$ より $\theta^2\leqq0.02$、$\theta\leqq\sqrt{0.02}\approx0.14$\,rad（約8°）}
\end{jibun}

\begin{reidai}[$x=a$ のまわりの展開]
$\sqrt{9.2}$ を、$a=9$ のまわりの2次近似で求めましょう。\par
\kangaekata ⑦の展開で、$f(x)=\sqrt{x}$、$a=9$、$x-a=0.2$ とする。$\sqrt9=3$。\par
\shiki $f(9)=3$、$f'(9)=\dfrac{1}{2\sqrt9}=\dfrac{1}{6}$、$f''(9)=-\dfrac{1}{4\times9\times\sqrt9}=-\dfrac{1}{108}$\par
\hspace*{2.6\zw}$3+\dfrac{1}{6}\times0.2-\dfrac{1}{108}\times\dfrac{0.2^2}{2}=3+0.033333-0.000185$\par
\kotae 3.033148（電卓3.033150）
\end{reidai}

\begin{jibun}
$\sqrt{16.4}$ を、$a=16$ のまわりの1次近似と2次近似で求めましょう。
\bkjans{$f'(16)=\frac{1}{8}$、$f''(16)=-\frac{1}{256}$。1次：$4+0.05=4.05$。2次：$4.05-\dfrac{0.16}{512}=4.049688$（電卓4.049691）}
\end{jibun}

\begin{reidai}[展開に式を入れる]
（1）$\cos2x$\quad（2）$e^{-x^2}$ のテイラー展開を、0でない項を3つずつ書きましょう。\par
\kangaekata ⑤の $\cos u$、④の $e^u$ の展開の $u$ に、$2x$ や $-x^2$ をかっこごと入れる（2-13）。\par
\shiki（1）$\cos u=1-\dfrac{u^2}{2!}+\dfrac{u^4}{4!}-\cdots$ に $u=2x$。$\dfrac{(2x)^2}{2!}=\dfrac{4x^2}{2}=2x^2$、$\dfrac{(2x)^4}{4!}=\dfrac{16x^4}{24}=\dfrac{2}{3}x^4$\par
\hspace*{2.6\zw}（2）$e^u=1+u+\dfrac{u^2}{2!}+\cdots$ に $u=-x^2$。$(-x^2)^2=x^4$ なので $\dfrac{(-x^2)^2}{2}=\dfrac{x^4}{2}$\par
\kotae （1）$1-2x^2+\dfrac{2}{3}x^4-\cdots$\quad（2）$1-x^2+\dfrac{x^4}{2}-\cdots$
\end{reidai}

\begin{jibun}
$\sin3x$ と $e^{2x^2}$ のテイラー展開を、0でない項を3つずつ書きましょう。
\bkjans{$\sin3x=3x-\dfrac{9}{2}x^3+\dfrac{81}{40}x^5-\cdots$。$e^{2x^2}=1+2x^2+2x^4+\cdots$}
\end{jibun}

\begin{reidai}[$\ln(1+x)$ の展開]
$\ln(1+x)$ を3次の項まで展開し、$\ln1.1$ を計算しましょう。\par
\kangaekata 3回微分して $x=0$ を入れ、③の式に入れる（5-14の $(\ln x)'=\frac{1}{x}$ と、合成関数の微分）。\par
\shiki $f'=\dfrac{1}{1+x}=(1+x)^{-1}$、$f''=-(1+x)^{-2}$、$f'''=2(1+x)^{-3}$\par
\hspace*{2.6\zw}$x=0$ で $f=\ln1=0$、$f'=1$、$f''=-1$、$f'''=2$ なので\par
\hspace*{2.6\zw}$\ln(1+x)\approx0+x+\dfrac{-1}{2!}x^2+\dfrac{2}{3!}x^3=x-\dfrac{x^2}{2}+\dfrac{x^3}{3}$\par
\hspace*{2.6\zw}$x=0.1$ で $0.1-0.005+0.000333$\par
\kotae $x-\dfrac{x^2}{2}+\dfrac{x^3}{3}$、$\ln1.1\approx0.09533$（電卓0.09531）
\end{reidai}

\begin{jibun}
$\ln1.2$ を、$\ln(1+x)$ の展開の3項で計算しましょう。
\bkjans{$0.2-0.02+0.002667=0.18267$（電卓0.18232）}
\end{jibun}

\begin{reidai}[次の項までたして、近似をよくする]
$\sqrt{1.02}$ を、$(1+x)^r\approx1+rx$ と、次の項 $\dfrac{r(r-1)}{2}x^2$ までたした式で計算しましょう。\par
\kangaekata $\sqrt{1.02}=(1+0.02)^{\frac{1}{2}}$ と見て、⑥の表の式を使う。\par
\shiki $r=\dfrac{1}{2}$、$x=0.02$。1次は $1+\dfrac{1}{2}\times0.02=1.01$\par
\hspace*{2.6\zw}次の項は $\dfrac{\frac{1}{2}\times(-\frac{1}{2})}{2}\times0.02^2=-\dfrac{1}{8}\times0.0004=-0.00005$\par
\kotae 1次で1.01、次の項までで1.00995（電卓1.0099505）
\end{reidai}

\begin{jibun}
$1.01^{20}$ を、$1+rx$ と、次の項までたした式で計算しましょう。
\bkjans{$1+20\times0.01=1.2$。次の項 $\dfrac{20\times19}{2}\times0.0001=0.019$ で1.219（電卓1.22019）}
\end{jibun}

\begin{reidai}[近似を式に入れる：振り子の高さ]
長さ\qty{2}{m}の振り子が、$\theta=0.2$\,rad ふれました。おもりは、いちばん下から何cm上がっていますか。$\cos\theta\approx1-\dfrac{\theta^2}{2}$ を使って求めましょう。\par
\kangaekata 支点からおもりまでの高さの差は $L\cos\theta$（直角三角形）。だから、いちばん下からの高さは $h=L-L\cos\theta=L(1-\cos\theta)$。ここに近似の式を入れる。\par
\shiki $1-\cos\theta\approx1-\Bigl(1-\dfrac{\theta^2}{2}\Bigr)=\dfrac{\theta^2}{2}$ なので $h\approx\dfrac{L\theta^2}{2}$\par
\hspace*{2.6\zw}$h\approx\dfrac{2\times0.2^2}{2}=0.04$（m）\par
\kotae 約\qty{4}{cm}（電卓で $2(1-\cos0.2)=0.0399$）
\end{reidai}

\begin{jibun}
長さ\qty{0.8}{m}の振り子が0.3\,rad ふれたときの、いちばん下からの高さを、同じ近似で求めましょう。
\bkjans{$h\approx\dfrac{0.8\times0.3^2}{2}=\qty{0.036}{m}$。約\qty{3.6}{cm}（電卓0.0357）}
\end{jibun}

%% ---------------- まちがえやすいところ・物理工学 ----------------
\bksection{tsumazuki}{まちがえやすいところをチェック}

\begin{tsumazuki}{$\sin\theta\approx\theta$ に度の値を入れる}
\bkng{$\sin10^\circ\approx10$}\bkok{$\sin10^\circ\approx0.1745$（$10^\circ=0.1745$\,rad）}
テイラー展開は、5-15の $(\sin x)'=\cos x$ から作ったので、ラジアンでしか成り立ちません。
\end{tsumazuki}

\begin{tsumazuki}{階乗でわるのをわすれる}
\bkng{$\dfrac{f'''(0)}{3}x^3$}\bkok{$\dfrac{f'''(0)}{3!}x^3=\dfrac{f'''(0)}{6}x^3$}
$x^3$ を3回微分すると $3\times2\times1=6$ が出てくるので、6でわります。3でわるのではありません。$\dfrac{f''(0)}{2}x^2$ の2も、ほんとうは $2!$ です。
\end{tsumazuki}

\begin{tsumazuki}{$x$ が大きいのに、少ない項で近似する}
テイラー展開は「$x=a$ の近く」で合う式です。$\sin3\approx3$ は大まちがい（正しくは0.14）。遠い所では、項をたくさんとるか、近い点 $a$ のまわりで展開しなおします。
\end{tsumazuki}

\begin{column}{振り子時計とガリレオ}
ガリレオは、教会のランプのゆれを脈ではかり、ゆれが小さくなっても周期が変わらないことに気づいたといわれます。のちにホイヘンスが振り子時計を作り、時計はとても正確になりました。ただし周期が変わらないのは、ふれが小さいとき（$\sin\theta\approx\theta$）だけです（⑧の図）。
\end{column}

\begin{butsuri}
\begin{itemize}[leftmargin=1.2\zw, itemsep=1pt, topsep=0pt, label={\color{bkPurple}●}]
\item \textbf{振り子とばね}\par\noindent 振り子は、ふれが小さいときだけ単振動。ばねも、のびが小さいときだけ $F=kx$ です。「小さいときは1次で近似」が、物理の式の多くを支えています。
\item \textbf{トランジスタと「小信号」}\par\noindent ダイオードやトランジスタでは、電流が電圧の指数関数（$e^{V/V_T}$ の形。$V_T$ は室温で約\qty{0.026}{V}）で変わります。電圧がほんの少しだけ変わるときは、$e^x\approx1+x$ のように1次で近似して、増幅の計算をします。
\item \textbf{ロボットと制御（9-10）}\par\noindent 手のひらに立てたほうき（倒立振子）の式にも $\sin\theta$ が入ります。立った姿勢のまわりで $\sin\theta\approx\theta$ と1次の式にし（線形化）、第9巻の方法で制御します。
\end{itemize}
\end{butsuri}

%% ---------------- 章末問題 ----------------
\bkendproblems
\bkproblemgroup{基本}{3}

\bkproblem 次の値を求めましょう。\quad（1）$4!$\quad（2）$\dfrac{6!}{4!}$\quad（3）$\dfrac{1}{5!}$ を小数で
\bkpans{（1）\nobreak 24\quad（2）\nobreak 30\quad（3）\nobreak 0.008333…}

\bkproblem $e^{3x}$ と $\ln(1+x)$ の、$x=0$ のまわりの2次近似を求めましょう。
\bkpans{$e^{3x}\approx1+3x+\dfrac{9}{2}x^2$。$\ln(1+x)$：$f'=\frac{1}{1+x}$、$f''=-\frac{1}{(1+x)^2}$ より \mbox{$x-\dfrac{x^2}{2}$}}

\bkproblem テイラー展開を使って、次の値を小数4けたまで求めましょう。\newline（1）$\sin0.1$（2項）\quad（2）$e^{0.2}$（4項）
\bkpans{（1）\nobreak $0.1-0.000167=0.0998$\quad（2）\nobreak $1+0.2+0.02+0.001333=1.2213$}

\bkproblem 長さ\qty{2.45}{m}の振り子の周期を求めましょう（$g=\qty{9.8}{m/s^2}$）。
\bkpans{\mbox{$T=2\pi\sqrt{0.25}=\pi\approx3.1$\,秒}}

\bkproblem $\sqrt{25.5}$ を、$a=25$ のまわりの1次近似で求めましょう。
\bkpans{$5+\dfrac{1}{10}\times0.5=5.05$（電卓5.04975）}

\bkproblem $\ln1.05$ を、$\ln(1+x)$ の展開の2項で計算しましょう。
\bkpans{$0.05-0.00125=0.04875$（電卓0.04879）}

\bkproblem 長さ\qty{0.5}{m}の振り子が小さくゆれるときの周期を求めましょう。ただし、$g=\qty{9.8}{m/s^2}$、$\sqrt{0.051}=0.226$ とします。
\bkpans{\mbox{$2\pi\sqrt{\dfrac{0.5}{9.8}}=2\times3.14\times0.226\approx1.4$\,秒}}

\bkproblemgroup{チャレンジ}{4}

\bkproblem $\cos\theta\approx1-\dfrac{\theta^2}{2}$ を使って、長さ $L$ の振り子が角 $\theta$ だけふれたときの、いちばん下からの高さ $h=L(1-\cos\theta)$ を近似しましょう。$L=\qty{1}{m}$、$\theta=0.1$\,rad のとき、何cmですか。
\bkpans{$h\approx\dfrac{L\theta^2}{2}$。$\dfrac{1\times0.01}{2}=\qty{0.005}{m}=\qty{0.5}{cm}$}

\bkproblem $(1+x)^r\approx1+rx$ を使って、$1.02^{10}$ と $\sqrt{1.06}$ の近似値を求めましょう。また、次の項 $\dfrac{r(r-1)}{2}x^2$ までたすと、$1.02^{10}$ はいくつになりますか（電卓では1.21899）。
\bkpans{$1.02^{10}\approx1.2$、$\sqrt{1.06}\approx1.03$。次の項 $45\times0.0004=0.018$ をたすと1.218}

\bkproblem $\cos x$ のテイラー展開の各項を微分すると、$-\sin x$ のテイラー展開になることを確かめましょう。また、$e^x$ の展開を項ごとに微分すると、どうなりますか。
\bkpans{$\Bigl(1-\dfrac{x^2}{2!}+\dfrac{x^4}{4!}-\cdots\Bigr)'=-x+\dfrac{x^3}{3!}-\cdots=-\Bigl(x-\dfrac{x^3}{3!}+\cdots\Bigr)$（$\frac{n}{n!}=\frac{1}{(n-1)!}$）。$e^x$ は、微分してももとの展開と同じ（$(e^x)'=e^x$）}

\bkproblem $\sin1$ を、テイラー展開の2項と3項で計算し、電卓の値0.841471とくらべましょう。ずれが0.001より小さくなるのは、何項からですか。
\bkpans{2項：$1-\dfrac{1}{6}=0.8333$（ずれ0.008）。3項：$+\dfrac{1}{120}$ で0.84167（ずれ0.0002）。3項から}

\begin{furikaeri}
\begin{checklist}
\item 2次近似で、曲がり方まで合わせる意味を説明できる。\hfill{\footnotesize\sffamily\color{bkTheme}→\pageref{bk:5-17:niji}ページ}
\item $e^x$・$\sin x$・$\cos x$ の展開を書き、値を計算できる。\hfill{\footnotesize\sffamily\color{bkTheme}→\pageref{bk:5-17:sin}ページ}
\item $\sin\theta\approx\theta$ から、振り子の周期を求められる。\hfill{\footnotesize\sffamily\color{bkTheme}→\pageref{bk:5-17:furiko}ページ}
\end{checklist}
\end{furikaeri}

\bkfukidashi{sora}{第5巻の道具が、これで全部そろった。最後に、表にまとめておこう！}
